\documentclass[11pt,reqno]{amsart}

%\usepackage{stix}
% -----------------------------------------------------------------------
\usepackage[english]{babel}
\usepackage{color}
\usepackage{amsmath, latexsym, amsthm, amsfonts,bm,amssymb,mathscinet} 
\usepackage{natbib}
\usepackage[text={5.8in,8.3in},centering]{geometry}
\usepackage{epsfig}
\usepackage[colorlinks,linkcolor=blue, citecolor=blue,urlcolor=blue]{hyperref}
\usepackage[dvipsnames]{xcolor}
\usepackage{mathtools}
\usepackage{nicefrac}
% -----------------------------------------------------------------------
\pagestyle{plain}

% -----------------------------------------------------------------------
\newcommand{\dd}{{\,\mathrm d}}
\newcommand{\si}{\sigma}
\renewcommand{\a}{\alpha}
\renewcommand{\b}{\beta}
\renewcommand{\th}{\theta}
\newcommand{\la}{\lambda}
\newcommand{\ga}{\gamma}
\newcommand{\ka}{\kappa}
\newcommand{\legeregel}{\par\vspace{1ex}\noindent}
\newcommand{\isd}{\stackrel{d}{=}}
\newcommand\Sgn{\mathop{\mathrm{Sgn}}\nolimits}
\newcommand{\eps}{\varepsilon}
\newcommand{\nphi}{\phi}
\renewcommand{\phi}{\varphi}
\newcommand{\asnaar}{\stackrel{\mathrm{a.s.}}{\longrightarrow}}
\newcommand{\pto}{\stackrel{\mathrm{p}}{\longrightarrow}}
\newcommand{\wto}{\rightsquigarrow}
\newcommand{\asto}{\stackrel{\mathrm{a.s.}}{\longrightarrow}}
\newcommand{\dto}{\stackrel{\mathrm{\scr{D}}}{\longrightarrow}}
\renewcommand{\scr}[1]{{\mathcal #1}}
\newcommand{\argmin}{\operatornamewithlimits{argmin}}
\newcommand{\argmax}{\operatornamewithlimits{argmax}}
\newcommand{\EE}{\mathbb{E}}
\newcommand{\FF}{\mathbb{F}}
\newcommand{\GG}{\mathbb{G}}
\newcommand{\PP}{\mathbb{P}}
\newcommand{\ind}{\mathbf{1}}
\newcommand{\RR}{\mathbb{R}}
\newcommand{\ZZ}{\mathbb{Z}}
\newcommand{\QQ}{\mathbb{Q}}
\newcommand{\CC}{\text{I\!\!\!C}}  
\newcommand{\NN}{\text{I\!N}}
\renewcommand{\vec}[1]{\underline{ #1}}
\newcommand{\ra}{\rightarrow}
\newcommand{\ua}{\uparrow}
\newcommand{\da}{\downarrow}
\renewcommand{\Pr}{\mathrm{Pr}}
\newcommand{\M}{\scr{M}}
\newcommand{\D}{\scr{D}}
\newcommand{\mb}[1]{\mathbf{ #1}}
\newcommand{\ft}[1]{\frametitle{#1}}
\newcommand{\Bm}{\begin{bmatrix}}
\newcommand{\Em}{\end{bmatrix}}
\newcommand{\bs}[1]{{\boldsymbol #1}}
\newcommand{\T}{{\prime}}

\DeclareMathOperator{\tr}{tr}
% -----------------------------------------------------------------------
% My definitions for theorems etc.
\theoremstyle{definition}
\newtheorem{thm}{Theorem}
\newtheorem{prop}[thm]{Proposition}
\newtheorem{lem}[thm]{Lemma}
\newtheorem{cor}[thm]{Corollary}
\newtheorem{rem}[thm]{Remark}
\newtheorem{ex}[thm]{Example}
\newtheorem{defn}[thm]{Definition}
\newtheorem{ass}[thm]{Assumption}
\newtheorem{exer}{Exercise}
% -----------------------------------------------------------------------
\newcommand{\pr}{\ensuremath{{\rm P}}}
\newcommand{\prob}[1]{\ensuremath{{\rm P}\!\left( #1 \right)}}
\newcommand{\prc}[2]{\ensuremath{{\rm P}( #1 \, |\, #2)}}
\newcommand{\expec}[1]{\ensuremath{{\rm E}\mspace{-1mu}\left[#1\right]}}
\newcommand{\var}[1]{\ensuremath{{\rm Var}\!\left( #1 \right)}}
\newcommand{\cov}[2]{\ensuremath{{\rm Cov}\!\left( #1 , #2 \right)}}
\newcommand{\corr}[2]{\ensuremath{{\rho}\mspace{-1mu}\left( #1 , #2 \right)}}
\newcommand{\mean}[1]{\ensuremath{\bar{#1}}}

\newcommand{\simind}{\stackrel{\rm ind}{\sim}}
\newcommand{\simiid}{\stackrel{\rm iid}{\sim}}
\newcommand{\pa}{\mathrm{pa}}
\newcommand{\KL}[2]{\mathrm{KL}\left(#1 \| #2\right)}

% Paper-specific definitions for TeX shortcuts
\def \ptil {\tilde{p}}
\def \pcirc {p^\circ}
\def \pstar {p^\star}
\def \Var {\text{Var}}
\def \R {\mathbb{R}}
\def \E {\mathbb{E}}
\def \Ncal {\mathcal{N}}

% g functions
\newcommand{\gfun}[3]{g_{#2}\left(#1_{#3}\right)}
\newcommand{\gfunx}[2]{g_{#1}\left(x_{#2}\right)}
% r functions
\newcommand{\rfun}[3]{r_{#2}\left(#1_{#3}\right)}
\newcommand{\rfunx}[2]{r_{#1}\left(x_{#2}\right)}
% U functions (nothing and *)
\newcommand{\Ufun}[3]{U\left(c_{#1}^{#2}, F_{#1}^{#2}, H_{#1}^{#2}\right)(#3)}
\newcommand{\Ufunx}[2]{U\left(c_{#1}^{#2}, F_{#1}^{#2}, H_{#1}^{#2}\right)(x)}
% U functions (bar)
\newcommand{\Ubar}[3]{U\left(\bar{c}_{#1}^{#2}, \bar{F}_{#1}^{#2}, \bar{H}_{#1}^{#2}\right)(#3)}
\newcommand{\Ubarx}[2]{U\left(\bar{c}_{#1}^{#2}, \bar{F}_{#1}^{#2}, \bar{H}_{#1}^{#2}\right)(x)}
% Kernels
\newcommand{\pullback}[2]{\left(\kappa_{#1, #2}, g_{#2}\right)(x)}
\newcommand{\pullbacktil}[2]{\left(\tilde{\kappa}_{#1, #2}, g_{#2}\right)(x)}


%% add notes via
\newcommand{\note}[1]{{\color{Red} #1}}

%---------------------------------------------------------------------------------------
\title{On the design of BFFG}
\author{Camilla \& Frank (VU Amsterdam)}
\date{\today}





%%%%%%%%%%%%%%%%%%%%%%5
\begin{document}
 \maketitle\tableofcontents

\begin{abstract}
Backward Filtering Forward Guiding is an algorithm for obtaining (weighted) samples of a partially observed stochastic process on a directed acyclic graph. For simplicity, we assume the state-space setting. The backward filtering step yields a sequence of {\it likelihood potentials} $\{g_i\}$ which are used in forward stochastic simulation of a {\it guided process} that is informed by these potentials. In this paper we investigate how properties of $\{g_i\}$ affect the efficiency of stochastic simulation of the guided process. 
\end{abstract}

\section{Introduction}

We consider a latent discrete-time Markov process $(x_1,\ldots, x_n)$. The edge from $0$ to $1$ fixes a prior on the initial state $x_1$.  At times $i=1,\ldots, n$ we assume an observation $v_i \mid x_i$ is generated from a distribution with density $k(x_i, v_i)$. We aim to sample from $x_{1:n} \mid v_{1:n}$. Let $h_i(x_i)$ be the density of $v_{i:n} \mid x_i$. The Markov process that is conditioned on $v_{1:n}$ is a Markov process with transition densities
\[ p^\star(x_i\mid x_{i-1})=  \frac{p(x_i\mid x_{i-1}) h_i(x_i)}{\int p(x_i\mid x_{i-1}) h_i(x_i)\dd x_i}.\]
 As the sequence of functions $\{h_i\}$ is often not available in closed form, we replace it with a sequence of {\it guiding} functions $\{g_i\}$. The resulting {\it guided} process has transition densities
\[ p^\circ(x_i\mid x_{i-1})=  \frac{p(x_i\mid x_{i-1}) g_i(x_i)}{\int p(x_i\mid x_{i-1}) g_i(x_i)\dd x_i}\]


The smoothing density is given by 
\begin{align*}
	p^\star(x_{0:n}) & = \prod_{i=0}^n p^\star(x_i\mid x_{i-1})= \prod_{i=0}^n \frac{p(x_i\mid x_{i-1}) h_i(x_i)}{\int p(x_i\mid x_{i-1}) h_i(x_i)\dd x_i}\\ &= \frac1{h_0(x_0)}\left[ \prod_{i=1}^n k(x_i, v_i)\right] \prod_{i=0}^n p(x_i \mid x_{i-1}), 
\end{align*} 
where $p(x_0 \mid x_{-1})\equiv p(x_0)$. The final equality is shown in Theorem 6 in \cite{bffg}. 
Under the guided process, we have 
\[ p^\circ(x_{0:n})= \prod_{i=0}^n p^\circ(x_i\mid x_{i-1}) = \prod_{i=0}^n \frac{p(x_i\mid x_{i-1}) g_i(x_i)}{\int p(x_i\mid x_{i-1}) g_i(x_i)\dd x_i}, \]
Therefore
\begin{equation}\label{eq:lr}
	 \frac{p^\star(x_{0:n})}{ p^\circ(x_{0:n})}=\frac1{h_0(x_0)} \Psi(x_{0:n}), 
	\end{equation}
	 where  \[ \Psi(x_{0:n})=\left[\prod_{i=1}^n k(x_i, v_i)\right]
 \prod_{i=0}^n \frac{\int p(x_i\mid x_{i-1}) g_i(x_i)\dd x_i}{g_i(x_i)}. \]


In particular, the likelihood is given by 
\begin{equation}\label{eq:likelihood}
	h_0(x_0) = \EE^\circ  \Psi(x_{0:n}). 
\end{equation} 



\subsection{Research questions / goals}
\begin{enumerate}
	\item Bound $\KL{\PP^\star}{\PP^\circ}$ by a discrepancy measure between $\{g_i\}$ and $\{h_i\}$. The behaviour of this KL-divergence characterises the quality of $\PP^\circ$ as importance sampling distribution to sample from $\PP^\star$ (Cf.\ \cite{chatterjee2018sample}).  
	\item The right-hand-side of \eqref{eq:likelihood} can be approximated by Monte-Carlo simulation from samples of the guided process.  Under which conditions is this estimator of finite variance? If we have a candidate $\{g_i\}$  which is not guaranteed to have this property, can we fix this by transforming each $g_i$?
	\item As we can view the guided process as a proposal in a guided Feynman-Kac model, can we use general results from particle filtering to study the quality of the guided process?
\end{enumerate}


\section{Formulation as guided Feynman-Kac model} 
In the notation of Chapter 5 in \cite{smcbook}, we have a guided Feynman-Kac model with 
\begin{align*}
\mathbb{M}_0(\dd x_0)&= p^\circ(x_0) \dd x_0 \\
 M_t(x_{t-1}, \dd x_t) &= p^\circ(x_t \mid x_{t-1}) \dd x_t \\
 f_t(v_t \mid x_t) &= k(x_t, v_t) \\
 G_t(x_{t-1}, x_t) &= \frac{k(x_t, v_t) p(x_t \mid x_{t-1})}{p^\circ(x_t\mid x_{t-1})}=k(x_t, v_t) \frac{\int p(x_t\mid x_{t-1}) g_t(x_t) \dd x_t}{g_t(x_t)}. 
\end{align*}
In a generic particle filtering algorithm, Cf.\ Algorithm 10.1 in \cite{smcbook}, the unnormalised weight is given by 
\[ w_t^n = G_t(X_{t-1}^{A_n}, X_t^n). \]
For particle filtering, by Proposition 11.3 in \cite{smcbook}, if the potentials $G_t$ are all upper bounded, then there exist constants $c_t'$ such that 
\[ \EE \left[ \frac1{N} \sum_{n=1}^N W_t^n \phi(X_t^n) - \mathbb{Q}_t(\phi)\right)^2 \le c'_t \frac{\|\phi\|^2_\infty}{N}, \]
where $W_t^n = w_t^n/\sum_{m=1}^N w_t^m$. 

The numerator of the expression for $G_t(x_{t-1},x_t)$ is bounded provided $k$ and $g_t$ are, which is a reasonable assumption in applications. However, the denominator, $g_t(x_t)$, is typically not bounded away from zero. For that reason, we think of {\it transforming} the maps $\{g_t\}$ to ensure this. For example via
\begin{itemize}
	\item[(A)] mapping $g_t(x)$ to $g_t(x)+\epsilon_t$;
	\item[(B)] mapping $g_t(x)$ to $\sqrt{g_t(x)^2+\epsilon_t}$,
\end{itemize}
where $\{\eps_t\}$ is a sequence of strictly positive numbers bounded away from zero. 
While such transformations on $g_t$ are simple, for BFFG to be efficient, at the same time we need that
\begin{enumerate}
	\item sampling under $p^\circ$ is simple;
	\item evaluation of $\int p(x_t\mid x_{t-1}) g_t(x_t) \dd x_t$ is simple. \note{I think a positive unbiased estimator would suffice.}
\end{enumerate}
Option (A) ensures (1) as
\begin{align*}
	 p^\circ(x_t \mid x_{t-1}) &=\frac{p(x_t \mid x_{t-1})(g_t(x_t)+\epsilon_t)}{\int p(x_t \mid x_{t-1})(g_t(x_t)+\epsilon_t) \dd x_t}\\& = \lambda_t(x_{t-1})\frac{p(x_t\mid x_{t-1}) g_t(x_t)}{\int p(x_t\mid x_{t-1}) g_t(x_t) \dd x_t}  + (1-\lambda_t(x_{t-1})) p(x_t\mid x_{t-1}), 
\end{align*}
where 
\[\lambda_t(x_{t-1})= \frac{\int p(x_t\mid x_{t-1}) g_t(x_t) \dd x_t}{\int p(x_t\mid x_{t-1}) g_t(x_t) \dd x_t+\epsilon_t}.\]
Hence, the guided process is a mixture distribution of the original guided process and the unconditional process.  \note{use rejection sampling for option (B)?}


%%%%%%%%%%%%%%%%%%%%%%%%%%%%%%%%%%%%%%%%%%%%%%%%%%%5

\section{KL-bounds}

Let  $\mathbb{P}_T^*$ and $\mathbb{P}_T^\circ$ denote the laws of  $X_{0:T}^*$ and  $X_{0:T}^\circ$ respectively. 
According to results in \cite{chatterjee2018sample}, the quality of an importance sampling distribution, which is $\PP^\circ_T$ here, is dictated by $\text{KL}(\mathbb{P}_T^*, \mathbb{P}_T^\circ)$ given by
\[
 \KL{\mathbb{P}_T^*}{\mathbb{P}_T^\circ)} = \int p^*(x_{0:T}) \log \left(\frac{\pstar(x_{0:T})}{\pcirc(x_{0:T})}\right) dx_{0:T} \, .
\]
We would like to bound this by a discrepancy measure between $\{h_i\}$ and $\{g_i\}$. \footnote{We should keep in mind that the effect of $h_i$ or $g_i$ is invariant under multiplying with a positive constant.} The following lemma is adapted from Proposition 3.1 in \cite{ju2021sequential}. 
\begin{lem}\label{lem:kl} Define $\xi_t(x)=\log h_t(x)-\log g_t(x)$. Then 
\begin{align*}
	\log \left(\frac{\pstar(x_{0:T})}{\pcirc(x_{0:T})}\right)&\le  \sum_{t = 0}^T \xi_t(x_t)  -  
\sum_{t = 0}^T\int \pcirc(x_t \mid x_{t - 1}) \xi_t(x_t) dx_t 
\end{align*}
	\end{lem}
\begin{proof}
We have 
\begin{align}\label{logratiosums}
\log \left(\frac{\pstar(x_{0:T})}{\pcirc(x_{0:T})}\right) &= \log \left(\prod_{t = 0}^T \frac{\pstar(x_t \mid x_{t-1})}{\pcirc(x_t \mid x_{t-1})}\right) = \log \left(\prod_{t = 0}^T \frac{h_t(x_t)}{g_t(x_t)} \frac{\int p(x_t \mid x_{t - 1}) g_t(x_t) dx_t}{\int p(x_t \mid x_{t - 1}) h_t(x_t) dx_t}\right) \notag\\
&= \sum_{t = 0}^T \log\left(\frac{h_t(x_t)}{g_t(x_t)} \right) + \sum_{t = 0}^T \log\left( \frac{\int p(x_t \mid x_{t - 1}) g_t(x_t) dx_t}{\int p(x_t \mid x_{t - 1}) h_t(x_t) dx_t}\right) \, \notag
%&= \sum_{t = 0}^T \log\left(\frac{h_t(x_t)}{g_t(x_t)} \right) + 
%\log\left( \frac{\int p(x_0) g_0(x_0) dx_0}{\int p(x_0) h_0(x_0) dx_0}\right) +
%\sum_{t = 1}^T \log\left( \frac{\int p(x_t \mid x_{t - 1}) g_t(x_t) dx_t}{\int p(x_t \mid x_{t - 1}) h_t(x_t) dx_t}\right),
\end{align}
using the convention $\pstar(x_0 \mid x_{-1}) = p(x_0)$ (and analogously for $\pcirc$).
Therefore, each  summand in the second summation of satisfies:
\begin{align*}
\log & \left(\frac{\int p(x_t \mid x_{t - 1}) g_t(x_t) dx_t}{\int p(x_t \mid x_{t - 1}) h_t(x_t) dx_t}\right) = \frac{\int p(x_t \mid x_{t - 1}) g_t(x_t) dx_t}{\int p(x_t \mid x_{t - 1}) g_t(x_t) dx_t} \cdot \log\left( \frac{\int p(x_t \mid x_{t - 1}) g_t(x_t) dx_t}{\int p(x_t \mid x_{t - 1}) h_t(x_t) dx_t}\right) \\ &\leq \frac{1}{\int p(x_t \mid x_{t - 1}) g_t(x_t) dx_t} 
\int p(x_t \mid x_{t - 1}) g_t(x_t) \log\left( \frac{p(x_t \mid x_{t - 1}) g_t(x_t)}{p(x_t \mid x_{t - 1}) h_t(x_t)}\right) dx_t \\
&=  
\int \frac{p(x_t \mid x_{t - 1}) g_t(x_t)}{\int p(x_t \mid x_{t - 1}) g_t(x_t) dx_t} \log\left( \frac{g_t(x_t)}{ h_t(x_t)}\right) dx_t \\&= \int \pcirc(x_t \mid x_{t - 1}) \log\left( \frac{g_t(x_t)}{ h_t(x_t)}\right) dx_t \, ,
\end{align*}
where the inequality follows from the logsum-inequality given in Lemma \ref{lem:logsumineq}. 
\end{proof}





\begin{cor}
	If $c_t \le \xi_t(x) \le C_t$, then 
	\[ \text{KL}(\mathbb{P}_T^*, \mathbb{P}_T^\circ) \le \sum_{t=0}^T (C_t-c_t). \]
\end{cor}
The condition in this corollary is admittedly very strong. 

\begin{cor}
We have
\begin{equation}\label{eq:klbound}
	\KL{\mathbb{P}_T^*}{\mathbb{P}_T^\circ)} \le  \sum_{t = 0}^T \left(\EE^\star \xi_t(X_t)-\sup_{x_{t-1}}\frac{p^\star(x_{t-1})}{p^\circ(x_{t-1})}  \EE^\circ \left[\min(\xi_t(X_t),0) \right]  \right), 
\end{equation} 
where $p^\star(x_{t-1})$ and $p^\circ(x_{t-1})$ denote the marginal densities of $x_{t-1}$ under $\PP^\star_T$ and $\PP^\circ_T$ respectively.
\end{cor}
\begin{proof}
	We have
\[ \KL{\mathbb{P}_T^*}{\mathbb{P}_T^\circ)} \le \sum_{t = 0}^T \EE^\star \xi_t(x_t)  -  
\sum_{t = 0}^T\int \int  \pcirc(x_t \mid x_{t - 1}) \xi_t(x_t) \dd x_t p^\star(x_{t-1}) d x_{t-1}. \] 
The integral in the second term an be written as (divide and multiply by $p^\circ(x_{t-1})$ in the integral)
\[ C:=-	\int \int \frac{p^\star(x_{t-1})}{p^\circ(x_{t-1})} p^\circ(x_{t-1}, x_t) \xi_t(x_t) \dd x_t \dd x_{t-1}. \]


Let 
\[ A=\{x_t \colon \xi_t(x_t)\ge 0\}.\]
We have
\[ C \le \int_{x_{t-1}} \int_{x_t \in A^c} p^\circ(x_{t-1}, x_t) (-\xi_t(x_t)) \dd x_t \frac{p^\star(x_{t-1})}{p^\circ(x_{t-1})} \dd x_{t-1} \]
as the integral will be negative when $x_t \in A$. 
Using Fubini, This leads to
\begin{align*}
	C &\le \sup_{x_{t-1}}\frac{p^\star(x_{t-1})}{p^\circ(x_{t-1})} \int_{x_t \in A^c} p^\circ(x_t)  (-\xi_t(x_t)) \dd x_t\\ &=\sup_{x_{t-1}}\frac{p^\star(x_{t-1})}{p^\circ(x_{t-1})} \EE^\circ \left[-\xi_t(X_t) \ind_{\{-\xi_t(X_t)>0\}}\right]\\
	&= -\sup_{x_{t-1}}\frac{p^\star(x_{t-1})}{p^\circ(x_{t-1})}  \EE^\circ \left[\min(\xi_t(X_t),0) \right] . 
\end{align*} 

\note{In \cite{ju2021sequential} it seems that $A=\emptyset$ is assumed. In that case, the term $\EE^\circ \min(\xi_t(X_t),0)$ can be replaced by $\EE^\circ \xi_t(X_t)$. }

\end{proof}



In specific cases where $\PP^\star_T$ is tractable, we may obtain sharper bounds.
\begin{ex}
	Gaussian case...
\end{ex}


\section{Argument by Ju/Heng}
We start from Lemma \ref{lem:kl}.

By Fubini's theorem applied to the second term in Lemma \ref{lem:kl}, we get\footnote{This is the first equation in the note by Jeremy Heng.}
\begin{cor}
If we define $p^{\star\circ}(x_0)=p^\circ(x_0)$ and for $t\ge 1$
\[ p^{\star\circ}(x_t) = \int p^\star(x_{t-1}) p^\circ(x_t \mid x_{t-1}) 	\dd x_{t-1}, \]
then
\[	\text{KL}(\mathbb{P}_T^*, \mathbb{P}_T^\circ) \le \sum_{t=0}^T \left(\EE^\star \xi_t(X_t) - \EE^{\star\circ} \xi_t(X_t)\right)
\]
with $\EE^{\star\circ}$ denoting expectation under $p^{\star\circ}$. 
\end{cor}
Then Heng defines $e_t = - \EE^{\star\circ} \xi_t(X_t)$. In his particular application apparently $g_T=h_T$ which implies   $e_T=0$.
We can write
\[ e_t = \EE^{\star\circ} \log \frac{g_t(X_t)}{h_t(X_t)} =\EE^{\star\circ} \log \frac{g_t(X_t)}{\tilde\psi_t(X_t)}+\EE^{\star\circ} \log \frac{\tilde\psi_t(X_t)}{h_t(X_t)}, \]
where 
\[ \tilde\psi(x_t) = p(y_t \mid x_t) \int p(x_{t+1}\mid x_t) g_{t+1}(x_{t+1}) \dd x_{t+1}. \]
This is the BIF recursion, but with $p$ rather than $\tilde p$ in the pullback. 
We have, using the logsum inequality
\begin{align*}
\log 	\frac{\tilde\psi_t(x_t)}{h_t(x_t)} &= \log \frac{\int p(x_{t+1}\mid x_t) g_{t+1}(x_{t+1}) \dd x_{t+1}}{\int p(x_{t+1}\mid x_t) h_{t+1}(x_{t+1}) \dd x_{t+1}}\\  & \le 
\frac1{\int p(x_{t+1}\mid x_t) g_{t+1}(x_{t+1}) \dd x_{t+1}} {\int p(x_{t+1}\mid x_t) g_{t+1}(x_{t+1})
\log \frac{g_{t+1}(x_{t+1})}{h_{t+1}(x_{t+1})}
 \dd x_{t+1}}\\ & = \int p^\circ(x_{t+1}\mid x_t)  \log \frac{g_{t+1}(x_{t+1})}{h_{t+1}(x_{t+1})}\dd x_{t+1}
\end{align*}
This implies 
\begin{align*}
	\EE^{\star\circ} \log \frac{\tilde\psi(X_t)}{h_t(X_t)} & \le \int \int p^\circ(x_{t+1}\mid x_t)  \log \frac{g_{t+1}(x_{t+1})}{h_{t+1}(x_{t+1})} p^{\star\circ}(x_t) \dd x_{t+1}\dd x_t  \\ & \le 
	M_t \int \int p^\circ(x_{t+1}\mid x_t)  \log \frac{g_{t+1}(x_{t+1})}{h_{t+1}(x_{t+1})} p^{\star}(x_t)\dd x_{t+1}  \dd x_t\\ &=
	M_t \int p^{\star\circ}(x_{t+1})  \log \frac{g_{t+1}(x_{t+1})}{h_{t+1}(x_{t+1})} \dd x_{t+1}\\&= M_t e_{t+1},
\end{align*}
where 
\[M_t = \sup_{x_t} \frac{p^{\star\circ}(x_t)}{p^\star(x_t)}.\]
Hence we obtain 
\begin{equation}\label{eq:bound_et}
	e_t \le \EE^{\star\circ} \log \frac{g_t(X_t)}{h_t(X_t)} + M_t e_{t+1},
\end{equation} 
to which we may apply Lemma \ref{lem:gronwalltype}. 
We conclude 
\begin{cor}
\[ \text{KL}(\mathbb{P}_T^*, \mathbb{P}_T^\circ) \le \sum_{t=0}^T\left(\EE^\star \xi_t(X_t) + e_t \right)  \]
where $e_t$ satisfies the bound \eqref{eq:bound_et}. 	
\end{cor}




So what remains to be shown?
\begin{enumerate}
	\item Control of $M_t$ (seems hard).
	\item Bounding $\EE^\star \xi_t(X_t) = \EE^\star \log \frac{h_t(X_t)}{g_t(X_t)}$. Cf.\ Equation (40) in Ju. \note{I think here we have to repeat the whole argument with splitting up the term using $\tilde\psi$.}
	\item Bounding $\EE^{\star\circ} \log \frac{g_t(X_t)}{\tilde\psi_t(X_t)}$. Cf.\ arguments in Equations (86) and (87) in Ju (here, the logsum inequality is used).
\end{enumerate}

Regarding (3), note that
\begin{align*}
	\log \frac{g_t(x_t)}{\tilde\psi_t(x_t)} & = \log \frac{\int \tilde p(x_{t+1}\mid x_t) g_{t+1}(x_{t+1}) \dd x_{t+1}}{\int p(x_{t+1}\mid x_t) g_{t+1}(x_{t+1}) \dd x_{t+1}}\\  & \le \int \frac{\tilde p(x_{t+1}\mid x_t) g_{t+1}(x_{t+1})}{\tilde p(x_{t+1}\mid x_t) g_{t+1}(x_{t+1}) \dd x_{t+1}} \log \frac{\tilde p(x_{t+1}\mid x_t)}{p(x_{t+1}\mid x_t)} \dd x_{t+1}
\end{align*}
by the logsum inequality. This term does not involve $h$ and should therefore by tractable in specific cases/applications. 

\subsection{On the proof in appendix B2 in \cite{ju2021sequential}}

This is about upper bounding $p^\star(x_t)/p^\circ(x_t)$. They argue that it suffices to bound $p^\star(x_{0:T})/p^\circ(x_{0:T})$. \note{Why does this suffice?} 
By \eqref{eq:lr} this boils down to bounding each of the terms in $\Psi(x_{0:n})$. Their argument, translated to our notation, is that $h_t$ is bounded and then 
$\int p(x_t \mid x_{t-1} h_t(x_t \dd x_t$ is bounded. Then, it suffices to lower bound $\int p(x_t \mid x_{t-1} g_t(x_t \dd x_t$. Their proof for this is specific to their setting.  

\subsection{On the proof in appendix B3 in \cite{ju2021sequential}}
This is about further bounding the bound in \eqref{eq:klbound}. In our notation, these bounds look like
\begin{align*}
	\EE^\star \log \xi_t(X_t) &\le \sum_{k=t}^{T-1} c^\star_{k+1} \sqrt{\EE^\star \phi^\star(X_t)^2} \sqrt{\EE^\star \Delta(X_t)} \\ -\EE^\circ \log \xi_t(X_t) &\le \sum_{k=t}^{T-1} c_{k+1} \sqrt{\EE^\circ\phi(X_t)^2} \sqrt{\EE^\circ \Delta(X_t)},
\end{align*}
where $\phi$, $\phi^\star$ and $\Delta$ are discrepancy measures of parameters in the true and approximate backward filters. 
%%%%%%%%%%%%%%%%%%%%%%%%%


\section{A toy example: Stochastic Volatility Model}
Consider a state-space model where the observations are given by the random sequence
\begin{equation}\label{r_obs}
    R_i = \exp\left(\frac{X_i}{2}\right) U_i, \quad 1 \leq i \leq n,
\end{equation}
and the latent states by
\begin{equation}\label{latent}
    X_i = \omega + \psi X_{i - 1} + \eta W_i, \quad 1 \leq i \leq n,
\end{equation}
where $(U_i)_{i = 1}^n$, $(W_i)_{i = 1}^n$ are independent random sequences of i.i.d.\ standard normal variates and $X$ has initial distribution $X_0 \sim \Ncal\left(\frac{\omega}{1 - \psi}, \frac{\eta^2}{1 - \psi^2}\right)$. We impose the parameter constraints $\omega \in \R$, $-1 < \psi < 1$ and $\eta > 0$.


We consider the following transformed version of the observation equation \eqref{r_obs}
\begin{equation}\label{v_obs}
    V_i \coloneqq \log\left(R_i^2\right) = X_i + U'_i, \quad 1 \leq i \leq n,
\end{equation}
where $U'_i \coloneqq \log\left(U_i^2\right)$, $i = 1, \dots, n$.

\subsection{Backward Filtering}
The proposed model is simple in the sense that the latent process is linear Gaussian. 
If the distribution of $U_i'$ would be Gaussian, the backward information filter provides a recursion that allows exact computation of $\{h_i\}$. This is unfortunately not the case, though can  obtain guiding functions  $\{g_i\}$ by backward filtering under the assumption
\begin{equation}\label{approx_cond_obs}
    V_i \mid X_i\sim \Ncal\left(\mu + X_i, \sigma^2\right)
\end{equation}
 with $\mu \coloneqq \E\left(U'_i\right) \approx -1.27$ and $\sigma^2 \coloneqq \Var\left(U'_i\right) = \nicefrac{\pi^2}{2}$.
In that case, the derivations are simple. It turns out to be useful to introduce the notation
\begin{equation}\label{Ufun}
    \Ufunx{}{} = \exp\left(c + Fx - \frac{1}{2}Hx^2\right).
\end{equation}
It turns out that \eqref{approx_cond_obs} implies any $g_i$ can expressed in this form for a triplet $(c_i, F_i, H_i)$. 

\begin{itemize}
	\item {\it Leaf Messages} From Theorem 6.1 (ii) in \cite{bffg}, if $v_n$ is observed at a leaf, then $\gfunx{n}{} = \Ufunx{n}{}$ with $c_n = \log\left(\varphi(v_n; \mu, \sigma^2)\right)$, $F_n = (v_n - \mu) \sigma^{-2}$ and $H_n = \sigma^{-2}$.
	\item {\it Pullback Kernels} By statement (iii) of the same theorem, we have for the pullback $\pullbacktil{n - 1}{n} = \Ubarx{n}{}$ with 
 \begin{align*}
	 \bar{H}_n &= \psi^2 / C_n \\ \bar{F}_n &= \psi (F_n/H_n - \omega) C_n^{-1} \\ \bar{c}_n &= c_n - \log\left(\varphi^{\text{can}}(0; F_n, H_n)\right) + \log\left(\varphi(\omega; F_n/H_n, C_n)\right),
	\end{align*}
	where $	C_n = \eta^2 + H_n^{-1}$.
\end{itemize}


Here 
\[ k(x,v)= \exp\left(v-x\right) f_{\chi^2(1)}(\exp(v-x)). \]
This is bounded since for $u>0$, the density of $u f_{\chi^2(1)}(u)$ behaves like 
$u u^{-1/2} e^{-u/2}$ which is bounded.
%\subsubsection*{Parametrization of $g$ functions} At time index $n$, we have
%\begin{equation}\label{gn}
%    \gfunx{n}{n} = \Ufunx{n}{\star},
%\end{equation}
%where $c_n^\star = c_n = \log\left(\varphi(v_n; \mu, \sigma^2)\right)$, $F_n^\star = F_n = (v_n - \mu) \sigma^{-2}$ and $H_n^\star = H_n = \sigma^{-2}$.
%
%For $i = 2, \dots, n$, we have the recursion
%\begin{equation}\label{grec}
%    \gfunx{i - 1}{i - 1} = \Ufunx{i - 1}{\star},
%\end{equation}
%where $c_{i - 1}^\star = c_{i - 1} + \bar{c}_n$, $F_{i - 1}^\star = F_{i - 1} + \bar{F}_n$ and $H_{i - 1}^\star = H_{i - 1} + \bar{H}_n$.

%%%%%%%%%%%%%%%%%%%%%%%%%5
\appendix
\section{Some results}

\begin{lem}\label{lem:logsumineq}
Writing $a = \int a(x) dx$ and $b = \int b(x) dx$, the integral analogue of the log-sum inequality is
\[
 a \log\left(\frac{a}{b}\right) \leq \int a(x) \log \left(\frac{a(x)}{b(x)}\right) dx  \, ,
\]
	\end{lem}
Ssee e.g. its use in \cite{csiszar2004information}, Lemma 7.4.



\begin{lem}\label{lem:gronwalltype}
Let $(e_t)_{t=1}^T$, $(f_t)_{t=1}^T$, and $(M_t)_{t=1}^{T-1}$ be sequences of nonnegative real numbers such that
\[
e_t \;\le\; f_t + M_t e_{t+1}, \qquad t=1,\dots,T-1,
\]
with terminal condition $e_T = 0$. Then
\[
\sum_{t=1}^T e_t 
\;\le\; \sum_{k=1}^T \Biggl( \sum_{t=1}^k \prod_{s=t}^{k-1} M_s \Biggr) f_k,
\]
where the empty product (when $t=k$) is taken to be $1$.
\end{lem}

\begin{proof}
We proceed by expanding the recursion. For $t = T-1$ we have
\[
e_{T-1} \le f_{T-1} + M_{T-1} e_T = f_{T-1} + M_{T-1}\cdot 0 = f_{T-1}.
\]
Inductively, suppose the claim holds for $e_{t+1},\dots,e_T$. Then
\[
e_t \;\le\; f_t + M_t e_{t+1}.
\]
Expanding $e_{t+1}$ gives
\[
e_t \;\le\; f_t + M_t f_{t+1} + M_t M_{t+1} f_{t+2} + \cdots + \Bigl(\prod_{s=t}^{T-1} M_s\Bigr) f_T.
\]
Summing this inequality over $t=1,\dots,T$ yields
\[
\sum_{t=1}^T e_t \;\le\; \sum_{t=1}^T \sum_{k=t}^T 
\Bigl(\prod_{s=t}^{k-1} M_s\Bigr) f_k.
\]
Reversing the order of summation shows that each $f_k$ appears with coefficient
\[
\sum_{t=1}^k \prod_{s=t}^{k-1} M_s,
\]
giving the stated bound.
\end{proof}



\section{Feynman-Kac models}
\subsection{Classical Feynman-Kac models (on chains)}

In the classical setting, a Feynman-Kac model describes a sequence of random variables \((X_0, X_1, \dots, X_n)\) evolving according to:

\begin{itemize}
  \item Markov transitions: \(X_{k+1} \sim M_k(X_k, \cdot)\)
  \item Potential functions: \(G_k(x_{k-1},x_k)\)
\end{itemize}

The associated Feynman-Kac measure on paths is defined as:
\[
\mathbb{Q}_n(dx_{0:n}) = \frac{1}{Z_n} \left[ \prod_{k=0}^n G_k(x_{k-1},x_k) \right] \mathbb{M}(dx_{0:n})
\]
where \(\mathbb{M}(dx_{0:n})\) is the law of the Markov chain, and \(Z_n\) is the normalizing constant.

\subsection*{Feynman-Kac models on directed trees}

In a directed tree, the process evolves on a rooted tree \(T\), where:

\begin{itemize}
  \item Each node \(t \in V(T)\) corresponds to a random variable \(X_t\)
  \item Each node \(t\) (except the root) has a parent \(\pa(t)\), and a transition kernel \(M_t(x_{\pa(t)}, \cdot)\)
  \item Each node has a potential function \(G_t(x_{\pa(t)}, x_t)\)
\end{itemize}

The root \(r\) has an initial distribution: \(X_r \sim \mu_r\)

\subsection*{Definition}

The Feynman-Kac measure \(\mathbb{Q}_T\) on the state space \(\{X_t \colon t \in T\}\) is given by:
\[
\mathbb{Q}_T(\dd\mathbf{x}) = \frac{1}{Z_T} \left[ \prod_{t \in T} G_t(x_{\pa(t)}, x_t) \right] \mathbb{M}_T(\dd\mathbf{x})
\]
where \(\mathbb{M}_T\) is the law of the Markov process over the tree:
\[
\mathbb{M}_T(d\mathbf{x}) = \mu_r(\dd x_r) \prod_{t \neq r} M_t(x_{\pa(t)}, \dd x_t)
\]
and \(Z_T\) is the normalizing constant (partition function):
\[
Z_T = \int \left[ \prod_{t \in T} G_t(x_{\pa(t)}, x_t) \right] \mathbb{M}_T(\dd\mathbf{x})
\]

%%%%%%%%%%%%%%%%%%%%%%%%%%%%%%%%%%%%%%%%%%%%%%%%%%%%
\bibliographystyle{alpha}
\bibliography{literature.bib}

\end{document}
